\begin{center}
\section*{M}
\end{center}
\addcontentsline{toc}{section}{M}
\term{Milestone}
Una milestone è un punto di riferimento significativo o un obiettivo importante all'interno di un progetto, utilizzato per misurare il progresso e il raggiungimento di determinati traguardi. Le milestone sono solitamente associate a date specifiche o a completamenti di specifiche attività o fasi di un progetto. Servono come punti di controllo per valutare se il progetto sta procedendo secondo i piani e per identificare eventuali ritardi o problemi. Le milestone possono essere utilizzate anche per comunicare i progressi del progetto alle parti interessate e per stabiliare scadenze chiare e tangibili per il completamento delle attività.
\term{Modello a V}
Il Modello a V è una rappresentazione grafica del ciclo di vita di un progetto di sviluppo software. Esso mostra in maniera dettagliata le relazioni fra le fasi di analisi e progettazione, la produzione e il testing e ciò che permette di minimizzare i rischi di progetto e di migliorare e garantire la qualità del prodotto.
